\documentclass[a4paper,11pt]{article}


\usepackage{amssymb,amsmath}
\usepackage[frenchb]{babel}
\usepackage[T1]{fontenc}
\usepackage[utf8]{inputenc}
\usepackage{hhline}
\usepackage{enumerate}
\usepackage{a4wide}
\usepackage{graphicx}

\newcommand{\T}{\mathbb{T}}
\newcommand{\N}{\mathbb{N}}
\newcommand{\Z}{\mathbb{Z}}
\newcommand{\Q}{\mathbb{Q}}
\newcommand{\R}{\mathbb{R}}
\newcommand{\C}{\mathbb{C}}
\newcommand{\K}{\mathbb{K}}


\def\om{\omega}
\def\eps{\varepsilon}
\def\DR{\mathcal{C}^\infty_c(\R)}

\DeclareMathOperator{\supp}{supp}

% lignes epaisses
\def\ligne#1{\leaders\hrule height #1\linethickness \hfill}

%-------------------environnements-----------------------------
\newcount\numexo
\numexo=1

\newenvironment{exercice}{\noindent\textbf{Exercice \number\numexo.---}}{\global\advance\numexo by 1}
\newenvironment{pb}{\noindent\textbf{Probl\`eme}}


\frenchspacing



%\usepackage{lmodern}
%
\usepackage{scrextend}
%%%%%%%%%%%%%%%%%%%%%%%%%%%%%%

%
\usepackage{helvet}
\renewcommand{\familydefault}{\sfdefault}
%\usepackage{rotating}
%
\thispagestyle{empty}
%
%\usepackage{inconsolata}
%
%
%\newcommand{\bigsize}{\fontsize{16pt}{20pt}\selectfont}

\usepackage{setspace}
\onehalfspacing


%%%%%%%%%%%%%%%%%%%%%%%%%%%%%%%%%%%%%%%%%%%%%%%%%%%%%%
% quelques symboles utiles en mathematiques
%%%%%%%%%%%%%%%%%%%%%%%%%%%%%%%%%%%%%%%%%%%%%%%%%%%%%%


%\newcommand{\R}{\mathbb{R}}
%\newcommand{\PP}{\mathbb{P}}
%\newcommand{\E}{\mathbb{E}}
%\newcommand{\N}{\mathbb{N}}
%\newcommand{\Q}{\mathbb{Q}}
%\newcommand{\Z}{\mathbb{Z}}

\newcommand{\sh}{{\mathop{\rm sh}}}
\newcommand{\ch}{{\mathop{\rm ch}}}
%\newcommand{\th}{{\mathop{\rm th}}}%commande deja d馭inie
\newcommand{\var}{\mathop{\rm Var}}
\newcommand{\cov}{\mathop{\rm Cov}}

                



\oddsidemargin  -0.25cm
\evensidemargin -0.25cm
\textwidth      17cm
%\headheight     0.in
%\headheight     -.25in
\topmargin      -2.5cm
\textheight     27cm





%%%%%%%%%%%%%%%%%%%%%%%%%%%%%%%%%%%%%%%%%%%%%%%%%


% les lettres grecques
\def\a{\alpha}
\def\b{\beta}
\def\l{\lambda}
%\newcommand{\a}{\alpha}\newcommand{\b}{\beta}\newcommand{\l}{\lambda}
%deja definies ailleurs
\newcommand{\g}{\gamma}
\newcommand{\m}{\mu}
\newcommand{\G}{{\Gamma}}



%%%%%%%%%%%%%% pour les lettres calligraphi仔s

\newcommand{\Arond}{{\mathcal A}}
\newcommand{\Brond}{{\mathcal B}}
\newcommand{\Crond}{{\mathcal C}}
\newcommand{\Drond}{{\mathcal D}}
\newcommand{\Erond}{{\mathcal E}}
\newcommand{\Frond}{{\mathcal F}}
\newcommand{\Grond}{{\mathcal G}}
\newcommand{\Hrond}{{\mathcal H}}
\newcommand{\Irond}{{\mathcal I}}
\newcommand{\Jrond}{{\mathcal J}}
\newcommand{\Krond}{{\mathcal K}}
\newcommand{\Lrond}{{\mathcal L}}
\newcommand{\Mrond}{{\mathcal M}}
\newcommand{\Nrond}{{\mathcal N}}
\newcommand{\Orond}{{\mathcal O}}
\newcommand{\Prond}{{\mathcal P}}
\newcommand{\Qrond}{{\mathcal Q}}
\newcommand{\Rrond}{{\mathcal R}}
\newcommand{\Srond}{{\mathcal S}}
\newcommand{\Trond}{{\mathcal T}}
\newcommand{\Urond}{{\mathcal U}}
\newcommand{\Vrond}{{\mathcal V}}
\newcommand{\Wrond}{{\mathcal W}}
\newcommand{\Xrond}{{\mathcal X}}
\newcommand{\Yrond}{{\mathcal Y}}
\newcommand{\Zrond}{{\mathcal Z}}












 \def \qed {{\hspace*{\fill}$\halmos$\medskip}}
 \def \Z {{\mathbb Z}}
 \def \R {{\mathbb R}}
 \def \C {{\mathbb C}}
 \def \N {{\mathbb N}}
 \def \P {{\mathbb P}}
 \def \E {{\mathbb E}}
 \def \Q {{\mathbb Q}}
 \def \ra {\rightarrow}
 \def \da {\downarrow}
 \def \ua {\uparrow}

 \def \cO {{\mathcal O}}
 \def \cW {{\mathcal W}}
 \def \cD {{\mathcal D}}
 \def \cL {{\mathcal L}}
 \def \cI {{\mathcal I}}
 \def \cJ {{\mathcal J}}
 \def \cR {{\mathcal R}}
 \def \cA {{\mathcal A}}
 \def \cM {{\mathcal M}}
 \def \cN {{\mathcal N}}
 \def \cE {{\mathcal E}}
 \def \cP {{\mathcal P}}
 \def \cQ {{\mathcal Q}}
 \def \cB {{\mathcal B}}
 \def \cF {{\mathcal F}}
 \def \cG {{\mathcal G}}
 \def \cC {{\mathcal C}}
 \def \cV {{\mathcal V}}
 \def \cH {{\mathcal H}}
 \def \cT {{\mathcal T}}
 \def \cZ {{\mathcal Z}}
 \def \cS {{\mathcal S}}


\begin{document}

 \begin{titlepage}
\newcommand{\HRule}{\rule{\linewidth}{0.5mm}}
%\text{}\\[2cm]
\text{}
\vspace{2cm}
\center

%\\[1cm]
\HRule \\[0.15cm]
{ \huge \bfseries Problèmes MATh.en.JEANS 2026-2027\\[0.15cm] }
\HRule \\[1.5cm]


%\\[1cm]
\vspace{2cm}
Lycée Honoré d'Estienne d'Orves - Carquefou\\
Lycée Grand Air - La Baule
\vspace{5cm}
\section*{Introduction}

C'est la soirée du siècle. 


\end{titlepage}

\newpage

\section{Un mélange parfait}

On souhaite mélanger un paquet de cartes. Comment s'y prendre ?


\newpage

\section{Casse-tête en solitaire }


Etude du Rubik's cube par la théorie des groupes.

%%Remarque :  on peut aussi réfléchir au problème avec des polygones flexibles, mais encore + dur .
\newpage

\section{Billard Nantais}

On se pose une question random sur un billard gênant.(compte comme le problème de géométrie ?)



\newpage

\section{}

Problème sur un karaoké ? (Dénombrement de chansons ?)

\newpage

\section{}

Problème de théorie des graphes ? Ou de géométrie. Ou autre, à l'image du problème 5
de l'année dernière.
Idées: point de fermat
       points entiers dans les cercles.
       


% \newpage


% \section{}

% \textbf{Défi 1}\\

% La traversée du désert (toutes les variantes de wikipédia possibles + complete refilling)\\

% \textbf{Défi 2}\\

% Le jeu du taquin : configurations possibles ou non, équivalence, que se passe-t-il lorsqu'on change la forme du taquin (configuration "garage", cf edouard lucas ) ?\\


% \textbf{Défi 3}\\

% Problème isopérimétrique : polygonal, approche d'une courbe par des polygones, preuve rigoureuse dans le cas triangulaire / polygonal ?


% \textbf{Défi 4}\\

% Des graphes et des lampes : 


% \textbf{Défi 5}\\

% Dénombrement nombres de catalan (photos de famille ?)



% \textbf{Défi 6}\\




% \textbf{Autres idées :}


\end{document}

%%% Local Variables:
%%% mode: latex
%%% TeX-master: t
%%% End:
